\documentclass[11pt]{article}
% Compile with: pdflatex draft.tex (run twice for cross-references)
%
% NOTE ON DOCUMENT CLASS: Automated Software Engineering (Springer) submissions ultimately require
% Springer Nature's official class (e.g. sn-jnl.cls, "sn-basic" / numbered-reference layout), which
% is not vendored in this repository. This draft uses the standard article class as a portable
% stand-in so it compiles anywhere; swap in Springer's class and its frontmatter commands when the
% official template is downloaded from the journal's submission site, mirroring how
% docs/research/jss/draft.tex uses elsarticle for the Elsevier-targeted companion submission.
%
% NOTE ON DRAFT STATUS (author-facing, not part of the manuscript): this draft supersedes the prior
% prescription-only manuscript, folding in the twenty-one-pattern anti-pattern catalog and its
% empirical validation (previously docs/antipatterns.md only) as a first-class Section 4, and
% re-mapping the three prescriptive operators (Section 6.3) explicitly onto the anti-patterns they
% repair. Two citation slots in Section 2.3 remain open (learning-based and LLM-based refactoring
% recommenders) and are marked as such rather than filled with invented references.

\usepackage{amsmath}
\usepackage{amssymb}
\usepackage{hyperref}
\usepackage{booktabs}
\usepackage{array}
\usepackage[margin=1in]{geometry}

\title{Graph-Based Detection of Architectural Anti-Patterns and Prescriptive Refactoring in
Distributed Publish--Subscribe Systems}
\author{Author Name(s) Here}
\date{}

\begin{document}
\maketitle

\begin{center}
\emph{Target venue: Automated Software Engineering (AuSE), Springer --- Special Issue on
``Intelligent techniques for CI/CD, DevOps, software evolution, technical debt analysis, and
refactoring recommendation.''}

\emph{SI topic mapping: this submission covers three of the bullet's five named areas directly ---
\textbf{technical debt analysis} (the anti-pattern catalog, \S4), \textbf{refactoring recommendation}
(the prescriptive operators, \S6), and \textbf{CI/CD/DevOps} (the delta-aware quality gate, \S7). No
AutoML/NAS/LLM contribution is claimed.}
\end{center}

\begin{abstract}
Distributed publish--subscribe middleware (ROS~2, DDS, MQTT) decouples producers and consumers, but
the resulting indirect dependency structure obscures how component failures cascade, and this
architectural technical debt accumulates invisibly to code-level static analysis (SCA) tools. Unlike
object-oriented design, where mature catalogs of named anti-patterns (God Class, Feature Envy,
Shotgun Surgery) give practitioners a shared vocabulary and testable detection rules, distributed
publish--subscribe architectures have no equivalent catalog: problems are discovered reactively,
through postmortems and cascade incidents, rather than proactively at design time. Even where
structural diagnostic frameworks exist, they typically operate \emph{open-loop}: they rank components
by criticality without naming the specific architectural pathology at fault or producing verified
guidance on how to repair it. We address both gaps with \textbf{SaG-Prescribe}, a graph-based
framework that (1) detects twenty-one named, severity-tiered publish--subscribe anti-patterns and bad
smells --- including Single Point of Failure, God Component, Broker Saturation (Hub-and-Spoke), Chatty
Pair, and QoS Policy Mismatch --- via formally specified topological signatures over thirteen
structural metrics and adaptive box-plot thresholds, each validated against independent
failure-simulation ground truth (Spearman $\rho = 0.876$ overall, rising to $\rho = 0.943$ at large
scale; $F_1 = 0.923$, precision $0.912$, recall $0.857$ for critical-tier classification); and (2)
compiles the resulting diagnosis --- components and patterns flagged \texttt{CRITICAL} or
\texttt{HIGH} --- into a transformation policy $\Delta(G)$ of three targeted graph-mutation operators
(logical topic splitting, physical anti-affinity reallocation, and transport QoS contract hardening),
each re-verified by the same discrete-event cascade simulator that produced the diagnosis before it
is surfaced as a recommendation. The pipeline is further operationalized as a delta-aware CI/CD
quality gate (\texttt{detect\_antipatterns.py}) that blocks only newly introduced structural
anti-patterns relative to a Git merge base, within sub-minute execution bounds.

Across eight benchmark scenarios used for detection validation and a seven-scenario subset used for
prescriptive evaluation --- spanning autonomous vehicles, IoT, finance, healthcare, and hyper-scale
enterprise topologies --- the generated prescriptions reduce the System Risk Index (SRI, lower is
better) in every scenario, with relative risk reductions between 1.4\% and 15.9\% (Wilcoxon
signed-rank across the seven baseline/mutated pairs: $W=0.0$, $p=0.0156$, significant at
$\alpha=0.05$). We report this result alongside a less flattering but more honest one: component-level
failure-impact reductions average only $+4.61\%$ and regress in two of seven scenarios ($-31.67\%$ and
$-25.36\%$), because the current prescriptive policy is applied greedily and unconditionally,
occasionally introducing new cascade hops that offset its own gains --- a finding that directly
motivates a per-edit acceptance filter as the highest-priority extension. The full generate--verify
loop completes in under 65 seconds for six of the seven prescription scenarios (4.7\,s--64.3\,s) and
in approximately 10.8 minutes for the 300-process enterprise topology, while the standalone CI/CD
detection gate completes in under 2 seconds to $\sim$40 seconds across the same scale range --- making
both proactive anti-pattern detection and closed-loop prescriptive refactoring practical as
merge-request-gated or nightly pre-deployment CI/CD stages.
\end{abstract}

\noindent\textbf{Keywords:} architectural anti-patterns; bad smells; technical debt analysis;
refactoring recommendation; publish--subscribe middleware; CI/CD quality gates; DevOps; failure
cascade simulation; graph mutation; search-based software engineering

% ===========================================================================
\section{Introduction}
% ===========================================================================

\subsection{Context and Motivation}

Distributed publish--subscribe middleware frameworks --- such as the Robot Operating System (ROS~2),
the Data Distribution Service (DDS), and MQTT --- form the communication backbone of modern
microservices, IoT systems, and safety-critical cyber-physical platforms. These architectures achieve
spatial, temporal, and synchronization decoupling among producers and consumers by routing messages
through shared topics and broker intermediaries. However, this decoupling introduces deep, non-linear
structural dependencies that obscure how component-level faults propagate through the wider system.
Hardening these networks against cascading failures requires proactive, continuous, pre-deployment
optimization before configurations are committed to runtime operational fabrics --- not remediation
after an incident.

\subsection{Two Open Gaps: An Unnamed-Pathology Gap and an Open-Loop Refactoring-Recommendation Gap}

Automated quality assurance has historically operated at the source-code level, through static code
analysis (SCA) platforms such as SonarQube. This produces an \textbf{Architecture-Code Gap}: a system
can have clean source code within every individual module yet remain highly fragile at the topology
level --- single points of failure (SPOFs), co-located deployment bottlenecks, or mismatched
communication attributes are invisible to file-scoped analysis. Shifting structural verification
``left'' into the CI/CD pipeline requires a paradigm shift from Static Code Analysis to \emph{Static
System Analysis}.

This gap is compounded by a second, narrower one that this paper addresses directly. Object-oriented
design has a mature vocabulary for structural pathology --- God Class, Feature Envy, Shotgun Surgery
--- each with a name, a formal detection rule, and an established refactoring strategy. Microservices
research has begun to build an analogous vocabulary for REST-based architectures (excessive
chattiness, shared databases, distributed monoliths). \textbf{No equivalent catalog exists for
publish--subscribe topologies.} Practitioners identify pub-sub-specific problems --- broker
saturation, topic fan-out explosion, QoS contract mismatches --- reactively, through postmortem
reports, performance regressions, or cascade failures, rather than proactively at design time, because
there is no shared name or testable rule for these conditions to check against.

Even where topology-aware diagnostic frameworks do exist, a further limitation remains. In our
companion paper \cite{SaG}, we introduced \textbf{Software-as-a-Graph (SaG)}, a static system analysis
framework that models pub-sub topologies as native heterogeneous graphs and produces diagnostic
criticality rankings $Q(v)$ and failure-impact predictions $I(v)$. SaG closes the Architecture-Code Gap
for diagnosis, but --- like refactoring-recommendation research more broadly, from code-smell
detectors to Search-Based Software Engineering (SBSE) --- it behaves as an \emph{open-loop diagnostic
engine}: it flags architectural vulnerabilities without synthesizing concrete, verified guidance on how
to resolve them, and it reports a numeric criticality score rather than naming the specific pathology
at fault. For architectural refactorings of distributed topologies, where the quality attribute is
resistance to cascading failure, this dual limitation is particularly hazardous: an architect cannot
act on an unnamed number, and an edit that looks beneficial locally can degrade global resilience if
its effect is never verified. We refer to the combination of these two limitations --- no named,
testable publish--subscribe anti-pattern vocabulary, and no verified guidance on how to repair what is
found --- as the \textbf{detection-and-remediation gap} this paper closes.

\subsection{Proposed Solution: SaG-Prescribe}

To close this gap, we present \textbf{SaG-Prescribe}, a graph-based framework unifying three stages:
detect, prescribe, and gate. First, SaG-Prescribe \textbf{detects} twenty-one named, severity-tiered
publish--subscribe anti-patterns and bad smells, each given a formal topological detection rule over
thirteen structural metrics and adaptive box-plot thresholds, and each validated against independent
failure-simulation ground truth. Second, it \textbf{prescribes}: components and patterns flagged
\texttt{CRITICAL} or \texttt{HIGH} --- itself informed by a \textbf{Code Quality Penalty (CQP)}
computed from static-analysis metrics --- feed a rule-based prescriptive engine that generates
targeted architectural mutations along three vectors: (1) logical topic splitting, isolating
high-fan-out publish channels; (2) physical anti-affinity reallocation, isolating co-located SPOFs;
and (3) transport QoS contract hardening, upgrading volatile/best-effort channels to reliable,
transient-local settings. Crucially, SaG-Prescribe implements \textbf{closed-loop simulation
verification}: each mutated candidate topology $G'$ is programmatically re-simulated with the same
discrete-event cascade simulator used for detection, under identical fault conditions, and a policy is
reported together with its verified change in the System Risk Index (SRI) rather than as an unverified
suggestion. Third, because the recommend-verify loop runs against an in-memory repository with no
database dependency, the same underlying engine is operationalized as a \textbf{delta-aware CI/CD
quality gate} that blocks only structural anti-patterns newly introduced relative to a Git merge base,
leaving pre-existing, risk-accepted debt untouched unless its severity worsens.

\subsection{Contributions}

\begin{enumerate}
\item \textbf{A catalog of twenty-one publish--subscribe anti-patterns and bad smells}, each with a
  formal, topology-based detection rule expressed over thirteen structural metrics and adaptive
  box-plot thresholds, organized into three severity tiers and four RMAV quality dimensions, and
  empirically validated against independent failure-simulation ground truth (\S4).
\item \textbf{A code-quality-augmented, closed-loop prescriptive refactoring pipeline} that translates
  the catalog's findings into counterfactual graph mutations via three named operators, each
  explicitly mapped to the anti-patterns it targets, and verifies each candidate against the same
  discrete-event cascade simulator that produced the diagnosis before it is surfaced as a
  recommendation (\S5--\S6).
\item \textbf{A DevOps-integrated, delta-aware CI/CD quality gate} (\texttt{detect\_antipatterns.py})
  that surfaces the catalog's twenty-one detectors as a blocking check, comparing candidate and
  merge-base topologies to flag only newly introduced anti-patterns, with a three-tier exit-code
  protocol and an auditable waiver register for accepted legacy risk (\S7).
\item \textbf{A multi-scenario empirical evaluation} across eight detection-validation scenarios and a
  seven-scenario prescriptive-evaluation subset, showing consistent, statistically significant SRI
  reductions of 1.4\%--15.9\% (Wilcoxon $W=0.0$, $p=0.0156$), an operator-level contribution analysis,
  a \textbf{component-level audit that reports mixed results honestly} rather than only the flattering
  system-level aggregate, and a runtime audit confirming compatibility with CI/CD budgets at every
  evaluated scale (\S8--\S9).
\end{enumerate}

The remainder of the paper is organized as follows. Section~2 surveys related work. Section~3
formalizes the graph model and the code-quality bridge. Section~4 presents the anti-pattern catalog
and its empirical validation. Section~5 defines the closed-loop optimization objective. Section~6
presents the prescriptive pipeline and its operators. Section~7 describes the DevOps/CI-CD
integration. Sections~8 and~9 present the experimental design and results. Section~10 discusses
implications and threats to validity, and Section~11 concludes.

% ===========================================================================
\section{Background and Related Work}
% ===========================================================================

\subsection{Publish--Subscribe Middleware Dependability}

Dependability research for message-oriented middleware historically centers on protocol verification,
fault-tolerant replication patterns, network traffic load balancing, and runtime contract validation.
Classical broker-fault-tolerance work replicates or partitions the broker itself to survive crashes
\cite{Pallickara2007,Chang2014}, and replicated-log designs such as Apache Kafka generalize this to a
durable, partitioned commit log underlying the pub-sub abstraction \cite{Wang2015}. Closer to our
DDS/ROS~2 setting, recent work analyzes the latency and reliability behavior of DDS's QoS-driven
retransmission protocol \cite{Lee2025Latency} and the static verifiability of interdependent DDS QoS
policies \cite{Lee2025QoS}. While Chaos Engineering practices inject faults into live clusters to
evaluate empirical resilience, this occurs late in the lifecycle and introduces operational risk.
SaG-Prescribe instead operates earlier, executing static system analysis on ``Architecture-as-Code''
descriptors to proactively detect and remediate structural pathology before deployment, treating the
topology as an open parameter rather than a fixed input.

\subsection{Anti-Pattern and Code-Smell Catalogs}

The most mature body of anti-pattern work addresses object-oriented design. Fowler's refactoring
catalog names recurring code-level pathologies (Long Method, Feature Envy, Shotgun Surgery) alongside
concrete refactoring transformations \cite{Fowler1999}. Brown et al.'s \emph{AntiPatterns} catalog
extends this template to the architectural and project-management level, formalizing a named pattern,
a recognizable symptom, and a remediation strategy as the standard specification unit
\cite{Brown1998}. Suryanarayana et al. systematize design smells with explicit, checkable structural
rules rather than purely qualitative description \cite{Suryanarayana2014}. Microservices research
builds an analogous vocabulary for REST-based architectures: Richardson catalogs recurring
microservices design and deployment patterns \cite{Richardson2018}, and Taibi et al. propose a taxonomy
of microservices-specific anti-patterns (excessive chattiness, shared databases, distributed
monoliths) grounded in practitioner surveys \cite{Taibi2020}.

The catalog presented in this paper (\S4) follows the same specification template --- a named pattern,
a formal detection rule, a remediation strategy --- but targets a domain none of the above cover: the
\emph{publish--subscribe communication topology} rather than object-oriented code structure or
request-response service boundaries. A pub-sub system can be architecturally pathological (a single
broker routing all traffic, a topic with hundreds of unmanaged subscribers) while every individual
component is internally well-structured by OO standards and every service boundary is REST-idiomatic.
The anomalies our catalog targets --- SPOF hosts, congested topic hubs, fragile transport contracts ---
have no file-level or service-boundary analog and are invisible to both OO and microservices catalogs.
Where those catalogs are grounded primarily in expert judgment and practitioner survey, \S4.5 further
validates each of our detection rules against independent failure-simulation ground truth.

\subsection{Refactoring Recommendation and Architectural Technical Debt}

Automated refactoring recommendation has been studied extensively at code level: smell-driven
recommenders detect structural anomalies (god classes, feature envy) and propose remediation
transformations, following the catalogs of \S2.2 \cite{Fowler1999,Suryanarayana2014}; learning-based
recommenders mine refactoring histories to predict refactoring opportunities [REF: ML-based
refactoring prediction]; and recent approaches employ large language models for refactoring suggestion
and explanation [REF: LLM-based refactoring]. At the architectural level, technical-debt research
quantifies the cost of structural decay and proposes prioritized remediation plans, in the spirit of
Brown et al.'s architectural AntiPatterns \cite{Brown1998}.

SaG-Prescribe differs from this body of work along two axes. First, its \textbf{scope} is the deployed
system topology --- applications, brokers, topics, hosts, and their QoS contracts --- rather than
source code within a module boundary; the anomalies it detects and remediates have no file-level
analog and are invisible to code-scope recommenders. Second, its \textbf{verification model} is
closed-loop: whereas code-level recommenders typically validate suggestions against static quality
metrics or historical acceptance data, SaG-Prescribe re-simulates every candidate topology against a
cascade failure model and surfaces only recommendations with verified risk improvements. In
technical-debt terms, the framework names architectural debt items (anti-patterns with quantified risk
via the Code Quality Penalty and RMAV attribution, \S3, and formal specification, \S4), proposes
repayments (mutation operators, \S6.3), and verifies the repayment's effect before recommending it ---
a verify-before-recommend discipline that, to our knowledge, has not been applied to pub-sub topology
refactoring.

\begin{quote}
\emph{[Citation-slot note, not part of the manuscript: the two remaining [REF: \ldots] markers above
must be populated from a real bibliography --- learning-based refactoring-prediction studies and
LLM-assisted refactoring work --- before submission. No references have been invented to fill these
slots.]}
\end{quote}

\subsection{Search-Based Software Engineering and Architecture Optimization}

Search-Based Software Engineering (SBSE) applies heuristic search to discover architectural
refactoring blueprints \cite{Harman2012}, and the architecture-optimization sub-field surveyed by
Aleti et al. \cite{Aleti2013} specifically targets quality-attribute-driven structural redesign.
However, classical search-based methods often operate open-loop, reporting recommendations without
verifying their operational efficacy against a cascade model. SaG-Prescribe combines the
multi-dimensional diagnostics of SaG \cite{SaG} with closed-loop simulation verification, ensuring
that every recommended edit is evaluated for its effect on the System Risk Index before acceptance.

\subsection{Diagnostic Foundation (SaG)}

We rely on the heterogeneous graph representation, multi-dimensional quality attribution (RMAV), and
discrete-event failure simulator of Software-as-a-Graph \cite{SaG} as our diagnostic baseline.
SaG-Prescribe builds directly on SaG's hexagonal ports, extending the domain service to name specific
structural pathologies (\S4) and to close the loop between diagnostic ranking and prescriptive
mutation (\S6). The full mathematical treatment of the diagnostic stages --- graph schema, projection
rules, RMAV attribution, and the learned failure-impact predictor --- is given in \cite{SaG} and is
not repeated here; \S3 summarizes only what the detection and prescriptive engines consume.

\subsection{Structural Criticality Analysis}

Graph-theoretic approaches offer constructs such as betweenness centrality \cite{Freeman1977},
PageRank, closeness, and articulation-point tests to pinpoint critical bridges; recent work applies
centrality measures directly to microservice dependency graphs to detect architectural anti-patterns
\cite{Bakhtin2025}, and complex-network analyses of software call graphs report the same small-world,
hub-dominated topologies that motivate criticality analysis in the first place \cite{Falci2017}.
Design-structure-matrix research on coupling and modularity \cite{Baldwin2000} and combinatorial
network-reliability theory \cite{Colbourn1987} provide the graph-theoretic foundations that the
anti-pattern catalog's structural signatures (\S4) build on; Lehman's laws of software evolution
\cite{Lehman1996} further motivate treating architectural criticality as a property that must be
re-checked continuously as a system evolves, rather than assessed once at initial design. Because
classical centrality metrics assume uniform edge semantics, they degrade on pub-sub layers, where
decoupled endpoints are separated by high-fan-out topics, brokers, and distinct QoS policies.
\cite{SaG} quantifies this gap and motivates the typed multigraph model that both the anti-pattern
catalog and SaG-Prescribe's mutation operators act on. These approaches identify fragility but, like
the diagnostic layer of SaG, stop short of naming the specific pathology and prescribing and verifying
remediation --- the step this paper automates.

% ===========================================================================
\section{System Model and Code-Quality-Augmented Technical Debt Analysis}
% ===========================================================================

\subsection{Heterogeneous Graph Formulation}

A distributed publish-subscribe deployment is modeled as a typed, weighted, directed multigraph
\begin{equation*}
G = (V, E, \tau_V, \tau_E, w_E, w_V)
\end{equation*}
where $\tau_V : V \to T_V$ partitions vertices into five semantic types,
\begin{equation*}
T_V = \{\text{Application}, \text{Library}, \text{Topic}, \text{Broker}, \text{Node}\}
\end{equation*}

\begin{itemize}
\item \textbf{Application} ($V_{\text{app}}$): active execution processes that produce or consume
  data.
\item \textbf{Library} ($V_{\text{lib}}$): shared code modules utilized across applications.
\item \textbf{Topic} ($V_{\text{topic}}$): named communication channels mediating message exchanges.
\item \textbf{Broker} ($V_{\text{broker}}$): middleware intermediaries routing message paths.
\item \textbf{Node} ($V_{\text{node}}$): physical or virtual hosting environments.
\end{itemize}

and $\tau_E : E \to T_E$ assigns each edge to a structural relation imported from the architecture
description,
\begin{equation*}
T_E = \{\text{PUBLISHES\_TO}, \text{SUBSCRIBES\_TO}, \text{ROUTES}, \text{RUNS\_ON}, \text{CONNECTS\_TO}, \text{USES}\}
\end{equation*}

\subsection{Derived Dependencies: The \texttt{DEPENDS\_ON} Projection}

To uncover logical dependency paths hidden behind decoupled pub-sub structures, the framework derives
explicit \texttt{DEPENDS\_ON} relations (directed \textbf{dependent $\to$ dependency}) via typed
projection rules:

\begin{itemize}
\item \textbf{Application-to-Application:} formed when a subscriber depends on a publisher via a
  shared topic.
\item \textbf{Application-to-Broker:} maps reliance on the specific broker instance routing an
  application's topics.
\item \textbf{Application-to-Library:} models the simultaneous blast radius where a shared library
  failure instantly impacts all consuming applications.
\item \textbf{Broker-to-Broker:} captures colocation vulnerabilities where multiple brokers share a
  physical host.
\end{itemize}

This projection produces $G_{\text{analysis}}$, organized across four architectural layers ---
\textbf{app} (applications only), \textbf{infra} (nodes only), \textbf{mw} (applications and
brokers), and \textbf{system} (all types) --- each providing a different lens for both the
anti-pattern detectors of \S4 and the RMAV attribution below.

\subsection{The Code Quality Penalty (CQP)}

To bridge local code quality with system architecture --- and to give this paper a direct, explicit
tie to the SI's ``software quality evaluation'' framing --- the framework ingests modular metrics from
static code analysis (SCA) APIs during model import. These features encompass total lines of code
(\texttt{cm\_total\_loc}), Weighted Methods per Class (\texttt{cm\_avg\_wmc}), Lack of Cohesion of
Methods (\texttt{cm\_avg\_lcom}), and the technical debt ratio (\texttt{sqale\_debt\_ratio}).
Rank-normalized properties map directly into a per-component \textbf{Code Quality Penalty}, defined
for Application and Library vertices:
\begin{equation*}
\mathrm{CQP}(v) = 0.10\,\text{loc\_norm} + 0.35\,\text{complexity\_norm} + 0.30\,\text{instability\_code} + 0.25\,\text{lcom\_norm}
\end{equation*}

CQP is the paper's single explicit channel from code-level quality signals into the architecture-level
risk model: it feeds directly into the Maintainability dimension of the RMAV attribution below, so a
module's static-analysis debt is not siloed from its topological criticality.

\subsection{Multi-Dimensional Quality Attribution (RMAV)}

Component criticality is decomposed into four orthogonal dimensions, ensuring that each structural and
code metric feeds exactly one perspective to preserve explanation legibility:

\begin{itemize}
\item \textbf{Reliability} ($R$): fault-propagation risk via Reverse PageRank (RPR) and fan-out
  concentration.
\item \textbf{Maintainability} ($M$): coupling complexity driven by betweenness centrality ($BT$),
  efferent QoS out-degree ($w\_out$), and the CQP metric:
\begin{equation*}
M(v) = 0.35\,\mathrm{BT}(v) + 0.30\,\mathrm{w\_out}(v) + 0.15\,\mathrm{CQP}(v) + 0.12\,\mathrm{CouplingRisk\_enh}(v) + 0.08\,(1-\mathrm{CC}(v))
\end{equation*}
\item \textbf{Availability} ($A$): single-point-of-failure risk via directed cut-vertex tests and
  QoS-amplified SPOF scores.
\item \textbf{Vulnerability} ($V$): exposure to adversarial reach, mapping attack propagation vectors.
\end{itemize}

These four profiles blend into a composite criticality score $Q(v)$ using Analytic Hierarchy Process
(AHP) weights \cite{Saaty1980} mixed with a uniform prior ($\lambda=0.70$) to prevent extreme parameter
concentration, yielding final weights of $(0.38, 0.24, 0.19, 0.19)$ for availability, reliability,
maintainability, and vulnerability respectively. Composite scores are mapped to five criticality tiers
(\texttt{CRITICAL}, \texttt{HIGH}, \texttt{MEDIUM}, \texttt{LOW}, \texttt{MINIMAL}) using adaptive
box-plot thresholding on the system's own score distribution (\texttt{CRITICAL}: $Q > Q_3 +
1.5\,\mathrm{IQR}$; \texttt{HIGH}: $Q_3 < Q \le$ upper fence). This section's typed graph, RMAV
dimensions, and adaptive box-plot machinery are the shared foundation consumed by both the anti-pattern
catalog (\S4) and the prescriptive engine (\S6).

% ===========================================================================
\section{A Catalog of Architectural Anti-Patterns for Publish--Subscribe Systems}
% ===========================================================================

\subsection{Anti-Patterns vs. Bad Smells}

Following the taxonomy established in object-oriented design research (\S2.2), this catalog
distinguishes between two categories of finding. An \textbf{anti-pattern} is a recognizable structural
configuration known to cause problems: it represents a deliberate or accidental architectural decision
that creates systemic risk and typically requires significant restructuring to resolve. A \textbf{bad
smell} is a surface symptom that suggests an underlying problem may exist --- not definitively harmful
in every context, but a reliable signal worth investigating, and often addressable with only a
localized change. In practice, the distinction is one of confidence: anti-patterns have
well-understood failure modes, whereas bad smells are heuristics that require human judgment to
confirm.

The key enabling insight is that architectural decisions in publish--subscribe systems leave
\textbf{measurable structural fingerprints} in the dependency graph: a single broker serving all
applications becomes an articulation point; a component that publishes to and subscribes from
everything exhibits extreme betweenness centrality; a topic with hundreds of subscribers shows
anomalous out-degree in the topic projection. Because these fingerprints are computable from the
system's static architecture --- the YAML configuration, the launch file, the infrastructure-as-code
--- without running the system at all, detection can occur proactively, at design time or during
CI/CD pipeline execution, before any deployment.

\subsection{Detection Methodology}

Detection operates over a thirteen-element metric vector $M(v)$ computed per component: PageRank
$\mathrm{PR}(v)$, Reverse PageRank $\mathrm{RPR}(v)$, betweenness centrality $\mathrm{BT}(v)$,
closeness centrality $\mathrm{CL}(v)$, eigenvector centrality $\mathrm{EV}(v)$, in- and out-degree
($\mathrm{DG}_{\text{in}}(v)$, $\mathrm{DG}_{\text{out}}(v)$), clustering coefficient
$\mathrm{CC}(v)$, articulation-point score $\mathrm{AP}_c(v)$, bridge ratio $\mathrm{BR}(v)$, and
QoS-weighted degree measures ($w(v)$, $w_{\text{in}}(v)$, $w_{\text{out}}(v)$). Topological metrics use
rank-based normalization by default, since they are typically highly skewed (a single hub-broker may
have betweenness 50$\times$ the median, which would compress all other values under min-max scaling);
linear properties (LOC, complexity, LCOM, hardware capacity) use min-max normalization, since their
absolute magnitude differences are meaningful.

A central design choice for robustness is the use of \textbf{adaptive box-plot thresholds} rather than
fixed global constants: for a metric vector $X$, the outlier fence is $Q_3 + 1.5 \times \mathrm{IQR}$,
and a component is flagged when its value exceeds this fence. This gives three properties important
for cross-system detection: \textbf{scale invariance} (a ``high'' betweenness score means something
different in a 10- versus a 300-component system, and the threshold adapts automatically);
\textbf{distribution awareness} (the threshold derives from the system's own metric distribution,
avoiding both over-flagging dense systems and under-flagging sparse ones); and \textbf{theoretical
grounding} (the $1.5 \times \mathrm{IQR}$ rule identifies genuine statistical outliers relative to a
component's peers, matching the definition of an anti-pattern as a structurally anomalous
configuration). Several patterns additionally target coupling imbalance directly, following Martin's
Instability metric \cite{Martin2003} enriched with topological path complexity
(\texttt{CouplingRisk\_enh}, \S3.4).

\subsection{Catalog Overview}

The twenty-one patterns are organized into three severity tiers --- \texttt{CRITICAL} (structural risk
requiring immediate architectural intervention; no production deployment should proceed without
addressing these), \texttt{HIGH} (significant risk materially degrading reliability, availability, or
maintainability; should be addressed in the current development cycle), and \texttt{MEDIUM}
(accumulated technical debt or localized risk; tracked for medium-term remediation) --- and mapped
onto the four RMAV dimensions of \S3.4. Table~\ref{tab:catalog} summarizes the full catalog; formal
detection rules and detailed remediation strategies for every pattern are given in the companion
technical reference (\texttt{docs/antipatterns.md}) and are cited rather than reproduced here in full.

\begin{table}[htbp]
\centering
\caption{Anti-pattern catalog summary.}
\label{tab:catalog}
\begin{tabular}{@{}p{2.6cm}p{1.5cm}p{3.0cm}p{6.0cm}@{}}
\toprule
Pattern & Severity & Primary RMAV Dimension & Detection Signal \\
\midrule
SPOF & CRITICAL & Availability & Articulation point, QoS-weighted SPOF score \\
SYSTEMIC\_RISK & CRITICAL & Reliability & Share of CRITICAL-tier components $> 20\%$ \\
CYCLE & HIGH & Architecture (cross-cutting) & Strongly connected component / self-loop \\
GOD\_COMPONENT & CRITICAL & Maintainability & Extreme betweenness $\wedge$ CRITICAL maintainability \\
BOTTLENECK\_EDGE & HIGH & Availability & Edge betweenness outlier \\
BROKER\_OVERLOAD & HIGH & Availability & Broker availability $\ge 2\times$ median, or sole broker \\
DEEP\_PIPELINE & HIGH & Reliability & Path length $\ge \max(5, P_{75})$ \\
TOPIC\_FANOUT & MEDIUM & Reliability & Topic subscriber out-degree outlier \\
CHATTY\_PAIR & MEDIUM & Maintainability & Bidirectional edge-weight product $> \tau$ \\
QOS\_MISMATCH & MEDIUM & Reliability & Publisher/subscriber QoS-weight gap $> \tau$ \\
ORPHANED\_TOPIC & MEDIUM & Maintainability & Zero in- or out-degree on structural graph \\
UNSTABLE\_INTERFACE & MEDIUM & Maintainability & High CouplingRisk\_enh $\wedge$ high $M(v)$ \\
BRIDGE\_EDGE & HIGH & Availability & Graph-theoretic bridge \\
FAILURE\_HUB & CRITICAL & Reliability & Reliability outlier $\wedge$ above-median out-degree \\
CONCENTRATION\_RISK & MEDIUM & Reliability & Top-3 PageRank share $> 0.5$ \\
HUB\_AND\_SPOKE & MEDIUM & Maintainability & Low clustering coefficient $\wedge$ degree $> 3$ \\
TARGET & CRITICAL & Vulnerability & Security-criticality tier $\ge$ CRITICAL \\
EXPOSURE & HIGH & Vulnerability & HIGH security tier $\wedge$ high closeness \\
CHAIN & MEDIUM & Architecture (cross-cutting) & Degree-bounded linear weakly-connected subgraph \\
ISOLATED & MEDIUM & Architecture (cross-cutting) & Zero total degree \\
COMPOUND\_RISK & CRITICAL & Architecture (cross-cutting) & Co-occurring SPOF + God/Hub/Failure-Hub finding \\
\bottomrule
\end{tabular}
\end{table}

\subsection{Representative Pattern Walkthroughs}

We highlight five patterns spanning severity tiers, RMAV dimensions, and detection technique
diversity.

\paragraph{SPOF (Single Point of Failure).} A component $v$ whose removal disconnects the graph ---
formally an articulation point. Detection combines a binary structural test with a continuous
\textbf{QoS-weighted SPOF severity} $\mathrm{QSPOF}(v) = \mathrm{AP}_c(v) \times w(v)$, so a flagged
SPOF is both structurally load-bearing and operationally significant. Unlike a performance bottleneck,
a SPOF produces a hard availability cliff: the system works completely until the SPOF fails, at which
point dependent functionality becomes entirely unavailable. Remediation centers on introducing
redundancy (replicated brokers, active-passive failover, stateless horizontally-scalable extraction for
application SPOFs) and circuit-breaker patterns to bound failover latency \cite{Nygard2018}.

\paragraph{GOD\_COMPONENT.} A component simultaneously exhibiting extreme betweenness centrality and
CRITICAL-tier maintainability ($\mathrm{BT}(v) > 0.30 \wedge \mathrm{Level}(M(v)) = \mathrm{CRITICAL}$).
It sits at a disproportionate share of shortest paths while also being the hardest component to change
safely, concentrating change-proneness, failure impact, and cognitive complexity simultaneously.
Remediation follows the Strangler Fig pattern: incrementally extracting cohesive publish/subscribe
responsibility subsets into new, purpose-built components while the original remains functional.

\paragraph{BROKER\_OVERLOAD (Hub-and-Spoke).} The pub-sub-specific instantiation of the classical
Hub-and-Spoke topology anti-pattern: a broker handling at least $2\times$ the median broker's routing
load (or the sole broker in the system, flagged unconditionally). One of our eight validation scenarios
deliberately encodes this anti-pattern with only two brokers serving seventy applications across twelve
nodes; both brokers are correctly classified \texttt{CRITICAL}, with broker failure-impact scores
exceeding 50\% of total system applications, confirming that the availability metric correctly
identifies broker-level overload independent of any single-application-level signal.

\paragraph{CHATTY\_PAIR.} A pair of application components maintaining a bidirectional, high-weight
dependency through separate topics in each direction: $(u \to v) \wedge (v \to u) \in
E_{\text{depends}}$ with $\mathrm{edge\_score}(u{\to}v) \times \mathrm{edge\_score}(v{\to}u) >
\tau_{\text{chatty}}$. This pattern detects \textbf{logical coupling masquerading as decoupling}: the
pub-sub layer gives the appearance of independence, but the communication pattern reveals that the
pair cannot be independently deployed, scaled, or reasoned about, and the coupling is distributed
across the broker rather than visible in code. Remediation introduces a mediator component or applies
event-carried state transfer, replacing the bidirectional conversational pattern with a
unidirectional, reactive one.

\paragraph{QOS\_MISMATCH.} A dependency edge $(u, v)$ where the publisher's QoS weight falls
substantially below the subscriber's expected guarantee level ($w_{\text{publisher}}(u) <
w_{\text{subscriber}}(v) - \tau_{\text{qos}}$). This pattern is unique to QoS-bearing middleware: it
detects a \textbf{silent connectivity failure} risk --- in DDS/ROS~2 systems, incompatible QoS policies
can prevent the endpoint match from being established at all, with no compile-time warning, while both
endpoints appear healthy in isolation. Remediation includes a QoS policy registry enforced in CI, or a
dedicated QoS-bridging relay component when the publisher's constraints (e.g., a hardware driver
limited to \texttt{BEST\_EFFORT}) cannot be upgraded directly.

\subsection{Empirical Validation of Detection}

Findings are validated empirically through the failure simulation pipeline: for each flagged
component, the simulated impact score $I(v)$ --- computed by exhaustive component removal and cascade
propagation --- provides independent evidence that a topological signature corresponds to real
structural risk. Table~\ref{tab:detection-validation} summarizes the headline validation metrics.

\begin{table}[htbp]
\centering
\caption{Detection validation metrics (overall).}
\label{tab:detection-validation}
\begin{tabular}{@{}lcc@{}}
\toprule
Metric & Target & Achieved \\
\midrule
Spearman $\rho$ ($Q$ vs. $I$) & $\ge 0.70$ & \textbf{0.876} \\
$F_1$-Score (critical classification) & $\ge 0.90$ & \textbf{0.923} \\
Precision & $\ge 0.85$ & \textbf{0.912} \\
Recall & $\ge 0.80$ & \textbf{0.857} \\
Top-5 Overlap & $\ge 0.70$ & \textbf{0.80} \\
\bottomrule
\end{tabular}
\end{table}

At large scale (150--300+ components), $\rho$ rises to \textbf{0.943}, indicating that prediction
accuracy improves as system scale increases --- precisely the regime where manual architectural
review becomes least practical. Pattern-specific evidence corroborates the aggregate: SPOF
classification achieves $F_1 > 0.95$ in application-layer analysis, and the deliberately-encoded
Hub-and-Spoke scenario (\S4.4) confirms BROKER\_OVERLOAD detection independent of the aggregate
correlation figures. A baseline comparison shows the composite $Q(v)$ score consistently outperforming
single-metric alternatives (betweenness centrality alone: $\rho = 0.75$); degree centrality alone
reaches $\rho = 0.95$ in synthetic graphs, a known artifact of topology generators that structurally
force high-degree hubs into SPOF positions, rather than evidence that degree alone suffices on
real-world heterogeneous topologies.

Detection validation uses an eight-scenario suite (01--08) spanning autonomous-vehicle, IoT,
financial, healthcare, deliberately-anti-pattern (Hub-and-Spoke), microservices, enterprise, and a
deterministic ``Tiny Regression'' smoke-test topology. Two scenarios play a distinguished role:
\textbf{Scenario 06} (sparse microservices mesh) is the primary \textbf{precision stress test} --- a
well-designed topology should produce few or no findings, confirming detectors do not over-flag
well-structured systems --- and \textbf{Scenario 07} (300+ components) is the primary
\textbf{scalability benchmark}, confirming detection algorithms scale gracefully to enterprise-scale
deployments. The seven-scenario subset used for prescriptive evaluation in \S8--\S9 excludes the
smoke-test fixture, which carries no domain-representative topology and would not contribute
meaningful signal to the prescriptive-efficacy questions asked there.

% ===========================================================================
\section{Closed-Loop Optimization Objective}
% ===========================================================================

The prescriptive task is to compute a transformation policy $\Delta$ producing a mutated topology
$G' = \Delta(G)$ that minimizes the aggregate failure-impact profile across system vertices, subject
to a modification budget:
\begin{equation*}
\min_{\Delta} \sum_{v \in V} I^*_{\Delta(G)}(v) \quad \text{subject to} \quad \mathrm{Cost}(\Delta) \le \mathcal{B}
\end{equation*}
where $I^*(v)$ denotes the simulated failure impact of component $v$. The candidate set for $\Delta$
is exactly the components and patterns flagged \texttt{CRITICAL} or \texttt{HIGH} by the catalog of
\S4. In the present implementation the modification budget is unconstrained ($\mathcal{B} = \infty$):
the engine emits every mutation whose triggering rule fires over this candidate set, and the aggregate
objective is tracked through the System Risk Index (SRI, \S8.4) computed by the verification stage. A
policy is \textbf{accepted} if and only if its verified improvement satisfies $\Delta\mathrm{SRI} > 0$
--- the whole-policy acceptance gate actually implemented in \texttt{PrescribeService}
(\texttt{saag/prescription/service.py}) and documented canonically in \texttt{docs/prescription.md}. A
stricter per-edit margin criterion, $\Delta A > \kappa\,\sigma_{\text{seed}}$, is discussed as future
work (\S11.2) but is not implemented; under the deterministic simulator configuration used throughout
this evaluation, $\sigma_{\text{seed}} = 0$ (empirically verified across seeds 42--46, \S8.3), so the
two criteria coincide for the present results. Budget-constrained policy search --- selecting the best
subset of mutations under an explicit cost model rather than firing every triggered rule --- is
likewise deferred to future work (\S11.2).

% ===========================================================================
\section{The SaG-Prescribe Prescriptive Pipeline}
% ===========================================================================

\subsection{Hexagonal Core Abstraction}

The system uses a decoupled hexagonal (ports-and-adapters) design separating domain orchestration from
persistence and communication infrastructure. Persistence services implement the
\texttt{IGraphRepository} port: production deployments run the Bolt-driven \texttt{Neo4jRepository},
while the verification loop and test suites use an isolated, thread-safe \texttt{MemoryRepository}
requiring no database instance. This substitution is what makes repeated counterfactual re-simulation
cheap enough for CI/CD integration (\S9.5).

\subsection{Pipeline Stages}

SaG-Prescribe extends the diagnostic pipeline of \cite{SaG} with a detect--generate--verify loop:

\begin{itemize}
\item \textbf{Stages 1--5: Diagnostic foundation and anti-pattern detection.} Ingest JSON/YAML
  topology descriptions, compute multi-layered topological metrics, attribute component criticality
  (\S3.4), run the twenty-one anti-pattern detectors of \S4 over the resulting metric vectors and RMAV
  scores, model failure cascades with a discrete-event simulator, and validate predictive alignment
  against simulation ground truth (\S4.5).
\item \textbf{Stage 6: Prescriptive recommendation generation (this paper).} The engine consumes
  components and patterns categorized \texttt{CRITICAL} or \texttt{HIGH} by Stage 5 and compiles a
  policy $\Delta(G)$ from the three operators of \S6.3. The compiled policy is applied to an exported
  copy of the topology, and the mutated model $G'$ is passed back into the analysis stage under
  identical fault scenarios and seeds to re-evaluate system risk (\S6.4).
\item \textbf{Stage 7: Review interface.} Baseline metrics, detected anti-patterns, prescribed
  modifications, and verification deltas are serialized for the Next.js dashboard (\textbf{SMART}),
  which renders side-by-side interactive topologies so architects can review and approve modifications
  before code commitment. Recommendations remain advisory: the human architect is the final authority
  (\S10.2).
\end{itemize}

\subsection{Refactoring Operators, Mapped to the Anti-Patterns They Target}

Each operator is formalized as a typed graph mutation rule triggered by specific anti-pattern findings
from \S4, making the detect-to-prescribe hand-off explicit rather than an unstated correspondence
between two independently-described stages.

\paragraph{Operator 1 --- Logical topic splitting} targets \texttt{TOPIC\_FANOUT} and topic-hub
contributions to \texttt{GOD\_COMPONENT}. For a flagged Topic $t$ with publisher set $P(t) = \{a : (a,
t) \in \text{PUBLISHES\_TO}\}$ and $|P(t)| > 1$, the operator replaces $t$ with dedicated sub-topics
$\{t_a : a \in P(t)\}$, rewiring each publisher to its own sub-topic and re-attaching subscriber edges
to the resulting set. This confines each data feed to its target subscribers, bounding the structural
blast radius of the original high-fan-out hub, and duplicates broker routing links accordingly.

\paragraph{Operator 2 --- Physical anti-affinity reallocation} targets \texttt{SPOF},
\texttt{BROKER\_OVERLOAD} (Hub-and-Spoke), and co-location contributions to \texttt{COMPOUND\_RISK}.
For a physical Node $n$ hosting multiple flagged components, the operator emits reallocation
constraints $(c, n_{\text{from}}, n_{\text{to}})$ moving each co-located component $c$ beyond the
first to an isolating host $n_{\text{to}}$, rewriting the corresponding \texttt{RUNS\_ON} edge and
duplicating \texttt{CONNECTS\_TO} links to preserve network reachability. The emitted constraints
correspond directly to container-orchestration anti-affinity scheduling rules.

\paragraph{Operator 3 --- Transport QoS contract hardening} targets \texttt{QOS\_MISMATCH} and any
\texttt{CRITICAL}/\texttt{HIGH} topic with a volatile transport configuration (\texttt{BEST\_EFFORT}
reliability, \texttt{VOLATILE} durability). The operator upgrades the contract to \texttt{RELIABLE}
reliability and \texttt{TRANSIENT} durability (raising transport priority where applicable), hardening
the channel against message loss during cascades.

\subsection{Closed-Loop Verification}

The verification engine executes the following loop, guaranteeing complete independence between
candidate generation and the validation path by separating graph structures from simulated metrics:

\begin{enumerate}
\item \textbf{Export:} serialize the source topology into a flat JSON schema.
\item \textbf{Mutate:} apply the compiled policy $\Delta(G)$ (splits, reallocations, upgrades)
  unconditionally to the exported structures.
\item \textbf{Sandbox isolation:} seed a temporary, thread-safe \texttt{MemoryRepository} with the
  mutated JSON and re-derive \texttt{DEPENDS\_ON} edges.
\item \textbf{Simulation oracle:} run the full analysis--simulation--validation suite on the sandbox
  model under identical fault scenarios and seeds.
\item \textbf{Resilience delta quantification:} compute $\Delta\mathrm{SRI} =
  \mathrm{SRI}_{\text{baseline}} - \mathrm{SRI}_{\text{mutated}}$ to report verified structural
  alignment.
\end{enumerate}

The core threat to structural predictors is circular leakage, where features inadvertently read data
from downstream labels. The framework mathematically avoids this via a strict \textbf{independence
guarantee}: all code metrics, RMAV calculations, and anti-pattern detection operate strictly on
$G_{\text{analysis}}$ (the derived projection layers), whereas ground-truth labels and SRI evaluations
are derived separately from raw $G_{\text{structural}}$ simulation waves. No simulation parameters ever
feed back into diagnostic metrics or candidate generation.

% ===========================================================================
\section{DevOps Integration and Delta-Aware CI/CD Gating}
% ===========================================================================

\subsection{Automated Code Review Architecture}

To continuously govern structural quality during rapid code evolution, the same underlying engine is
operationalized as a blocking check script (\texttt{detect\_antipatterns.py}) in continuous
integration and delivery (CI/CD) pipelines, surfacing the twenty-one detectors of \S4 directly as a CI
check. Whenever an engineer alters system structures or configures new messaging topology, the gate
parses the ``Architecture-as-Code'' descriptors and populates an in-memory graph view --- reusing the
\texttt{MemoryRepository} substrate that makes the prescriptive loop itself CI/CD-viable (\S6.1).

\subsection{Delta-Aware Regression Semantics}

Absolute quality gates that fail builds on any critical structural anti-pattern are unsustainable in
industrial software development, since real architectures frequently contain intentional,
risk-accepted debt (e.g. legacy unreplicated components). To resolve this, the gate uses
\textbf{delta-aware semantics}: it compares the pull-request candidate topology against the target
branch's merge-base topology, and isolates and flags only \emph{newly introduced} anti-patterns.
Pre-existing findings pass unless their severity worsens, and intentional anomalies can be bypassed via
an auditable, time-bound \textbf{waiver register} --- mirroring the ``Clean as You Code'' discipline
familiar from code-scope quality platforms, applied here at topology scope: the gate blocks new
architectural debt, while the prescriptive engine (\S6) proposes verified repayments of existing debt.

\subsection{Exit-Code Protocol}

The quality gate enforces automated code review by terminating with standardized exit codes that
command CI/CD pipeline workers:

\begin{itemize}
\item \textbf{Exit Code 0:} no new anti-patterns detected; build passes, deployment permitted.
\item \textbf{Exit Code 1:} new \texttt{MEDIUM}-severity anti-patterns or QoS warnings introduced;
  build passes with warnings compiled into the developer's code-review dashboard.
\item \textbf{Exit Code 2:} new, unwaived \texttt{CRITICAL} or \texttt{HIGH} severity anomalies (e.g.
  newly introduced un-replicated SPOFs or routing loops) discovered; \textbf{the build breaks and
  deployment is blocked}.
\end{itemize}

Because the recommend-verify loop of \S6 runs within CI time budgets (\S9.5), prescriptive repayment
proposals can be regenerated on every merge request, keeping the architectural-debt register
synchronized with the evolving topology rather than only gating new regressions.

% ===========================================================================
\section{Experimental Design}
% ===========================================================================

\subsection{Research Questions}

\begin{description}
\item[RQ1 (Detection efficacy and precision)] Does the anti-pattern catalog correlate with
  ground-truth failure impact, and does it avoid over-flagging well-structured systems?
\item[RQ2 (Prescriptive efficacy)] Does the closed-loop engine reduce the System Risk Index across
  heterogeneous scenarios, and are the reductions statistically significant?
\item[RQ3 (Operator contributions)] How do the individual refactoring operators contribute to the
  observed improvements across topological regimes?
\item[RQ4 (Component-level remediation efficacy)] Do system-level SRI improvements translate uniformly
  into component-level failure-impact reductions, or can the greedy, unconditional policy application
  backfire on individual components?
\item[RQ5 (Computational overhead and CI/CD feasibility)] What is the wall-clock execution time of the
  detection gate and of the full prescriptive pipeline as system scale grows --- and is either
  compatible with CI/CD budgets?
\end{description}

\subsection{Benchmark Scenarios}

Detection validation (\S4.5, RQ1) uses an eight-scenario suite (01--08), including a deterministic
``Tiny Regression'' smoke-test fixture used as a CI regression sanity check. Prescriptive evaluation
(RQ2--RQ5) uses the seven-scenario subset (01--07) of that suite, spanning realistic domain verticals
and scale presets; the smoke-test fixture is intentionally excluded from prescriptive evaluation since
it carries no domain-representative topology and would not contribute meaningful signal to questions
about prescriptive efficacy or CI/CD runtime at scale.

\begin{table}[htbp]
\centering
\caption{Scenario scale and topology summary (seven-scenario prescriptive-evaluation subset).}
\label{tab:scenarios}
\begin{tabular}{@{}lcccccc@{}}
\toprule
Scenario & Apps & Libs & Topics & Brokers & Nodes & Struct. Edges ($|E|$) \\
\midrule
S01 Autonomous Vehicle & 80 & 20 & 40 & 4 & 8 & 797 \\
S02 IoT Smart City & 200 & 10 & 80 & 6 & 30 & 1322 \\
S03 Financial Trading & 60 & 18 & 35 & 5 & 6 & 580 \\
S04 Healthcare & 50 & 12 & 25 & 3 & 8 & 400 \\
S05 Hub-and-Spoke & 70 & 25 & 30 & 2 & 12 & 797 \\
S06 Microservices Mesh & 90 & 30 & 45 & 6 & 15 & 680 \\
S07 Hyper-Scale Enterprise & 300 & 50 & 120 & 10 & 40 & 3245 \\
\bottomrule
\end{tabular}
\end{table}

This seven-scenario suite is a subset of the eight-scenario preset library used across companion SaG
materials \cite{SaG}; the eighth preset (an ATM system) is reserved there for external validation of
the diagnostic model against a non-synthetic reference topology and is not reused here, since this
paper's prescriptive evaluation targets the closed-loop pipeline rather than diagnostic external
validity. Replaying SaG-Prescribe on the ATM topology is future work (\S11.2).

\subsection{Experimental Protocol}

For each scenario, the prescriptive engine runs in-memory over the topology, measuring the SRI before
and after mutation together with the counts of the three operators. The discrete-event cascade
simulator runs in its default deterministic configuration (cascade propagation threshold 0.2, cascade
probability 1.0, non-Poisson event arrivals). We empirically verified this determinism: re-running the
full generate--verify loop under five candidate seeds (42--46) produced byte-identical SRI values in
every scenario ($\sigma_{\text{seed}} = 0$), so a single run per scenario is reported and paired
comparisons are not confounded by simulation stochasticity under the current configuration. For RQ2,
we test the baseline-vs-mutated SRI pairs across the seven scenarios with a Wilcoxon signed-rank test.

\subsection{Metrics}

\textbf{Detection metrics (RQ1).} Spearman rank correlation $\rho$ between the composite criticality
score $Q(v)$ and simulated impact $I(v)$; $F_1$-score, precision, and recall for \texttt{CRITICAL}-tier
classification against simulation ground truth; Top-5 overlap between predicted and simulated
rankings.

\textbf{System Risk Index (SRI, RQ2--RQ5).} The primary prescriptive outcome measure is a composite
risk index over the four RMAV health dimensions:
\begin{equation*}
\mathrm{SRI} = 0.25\,(1 - H_R) + 0.25\,(1 - H_M) + 0.25\,(1 - H_A) + 0.25\,(1 - H_V)
\end{equation*}
where $H_R, H_M, H_A, H_V$ are the system-level Reliability, Maintainability, Availability, and
Vulnerability health scores. Lower SRI indicates lower composite structural risk. We report
$\Delta\mathrm{SRI} = \mathrm{SRI}_{\text{baseline}} - \mathrm{SRI}_{\text{mutated}}$ (positive =
improvement), per-operator recommendation counts, and --- separately --- component-level
failure-impact deltas to answer RQ4.

% ===========================================================================
\section{Results}
% ===========================================================================

\subsection{Detection Efficacy and Precision (RQ1)}

The catalog's detection rules correlate strongly with independent failure-simulation ground truth
across the eight-scenario detection suite: Spearman $\rho(Q, I) = 0.876$ overall, rising to $\rho =
0.943$ at large scale (150--300+ components); $F_1 = 0.923$, precision $0.912$, recall $0.857$ for
\texttt{CRITICAL}-tier classification; and Top-5 overlap of $0.80$ (\S4.5, Table~\ref{tab:detection-validation}).
The precision stress test (Scenario 06, sparse microservices mesh) confirms detectors do not
over-flag well-structured topologies, and the scalability benchmark (Scenario 07, 300+ components)
confirms the metrics above hold at enterprise scale rather than degrading. Pattern-specific evidence
--- SPOF $F_1 > 0.95$; the deliberately-encoded Hub-and-Spoke scenario correctly classifying both
brokers \texttt{CRITICAL} with broker-failure impact exceeding 50\% of system applications ---
corroborates that these aggregate figures are not an artifact of averaging over easy and hard cases.

\subsection{Prescriptive Efficacy (RQ2)}

\begin{table}[htbp]
\centering
\caption{Prescriptive optimization results across scenarios.}
\label{tab:prescriptive-results}
\begin{tabular}{@{}lccccccc@{}}
\toprule
Scenario & Base SRI & Mut. SRI & $\Delta$SRI & Rel. Red. & Splits & Reallocs & Upgrades \\
\midrule
S01 Autonomous Vehicle & 0.3645 & 0.3535 & 0.0110 & 3.0\% & 35 & 121 & 0 \\
S02 IoT Smart City & 0.4206 & 0.3537 & 0.0669 & 15.9\% & 58 & 276 & 51 \\
S03 Financial Trading & 0.3675 & 0.3482 & 0.0193 & 5.3\% & 31 & 88 & 6 \\
S04 Healthcare & 0.3809 & 0.3757 & 0.0052 & 1.4\% & 19 & 74 & 6 \\
S05 Hub-and-Spoke & 0.3595 & 0.3527 & 0.0068 & 1.9\% & 30 & 97 & 0 \\
S06 Microservices Mesh & 0.3612 & 0.3542 & 0.0070 & 1.9\% & 40 & 123 & 0 \\
S07 Hyper-Scale Enterprise & 0.3614 & 0.3469 & 0.0145 & 4.0\% & 119 & 409 & 0 \\
\bottomrule
\end{tabular}
\end{table}

In all seven scenarios the mutated topology achieves a lower SRI than the baseline (mean
$\Delta$SRI~=~0.0187). A Wilcoxon signed-rank test over the seven baseline/mutated pairs yields $W =
0.0$, $p = 0.0156$ --- significant at $\alpha = 0.05$, since all seven pairs share the same sign of
difference (the smallest attainable two-sided $p$-value at $n=7$).

The largest improvement occurs in \textbf{S02 (IoT Smart City)}, where SRI drops by 0.0669 (15.9\%),
driven by 51 QoS upgrades stabilizing high-loss best-effort links (repairing \texttt{QOS\_MISMATCH}
findings, \S6.3) combined with anti-affinity constraints that partition key microservices. The
smallest improvement occurs in \textbf{S04 (Healthcare, 1.4\%)}: its centralized monitoring fan-outs
are intentional and already run under long durability horizons, leaving less structural headroom for
the current operator set. In \textbf{S05 (Hub-and-Spoke)}, topic splitting mitigates the deliberate
broker-pair \texttt{BROKER\_OVERLOAD} bottleneck (0.3595 $\to$ 0.3527).

\subsection{Operator Contributions (RQ3)}

Operator activity tracks each scenario's topological regime and the anti-patterns each scenario
stresses (\S6.3):

\begin{itemize}
\item \textbf{Logical topic splits} dominate where publisher fan-out concentrates: S07 receives 119
  splits, alleviating \texttt{TOPIC\_FANOUT} and \texttt{GOD\_COMPONENT} congestion in heavily shared
  publisher--subscriber hubs.
\item \textbf{Physical reallocations} are the most frequently emitted mutation overall (409 in S07,
  276 in S02), reflecting how often \texttt{SPOF} and co-location \texttt{COMPOUND\_RISK} findings
  arise in dense deployments; anti-affinity constraints resolve them without touching logical
  structure.
\item \textbf{QoS upgrades} activate almost exclusively in high-loss profiles --- 51 in S02 --- where
  hardening \texttt{QOS\_MISMATCH} volatile/best-effort channels yields the single largest SRI gain
  observed. Scenarios already running reliable contracts (S01, S05, S06, S07) receive none,
  confirming the rule does not fire spuriously.
\end{itemize}

An ablation applying each operator class in isolation would let us attribute $\Delta$SRI per operator
rather than inferring contribution from counts; we flag this as future work (\S11.2) rather than
including a speculative table here.

\subsection{Remediation Efficacy at Component Boundaries (RQ4)}

To thoroughly audit the refactoring engine beyond the system-level aggregate, we track individual
component-level failure impacts before and after applying mutations. Because the current
implementation applies the fully compiled policy unconditionally without a per-edit acceptance filter,
it delivers a mixed result that exposes the structural trade-offs of greedy graph mutation:

\begin{itemize}
\item \textbf{S02 (IoT Smart City):} $+47.01\%$ mean reduction in individual component failure impact.
\item \textbf{S04 (Healthcare):} $+30.94\%$ component-level risk improvement.
\item \textbf{S05 (Hub-and-Spoke):} $-31.67\%$ degradation in mean component metrics.
\item \textbf{S07 (Hyper-Scale Enterprise):} $-25.36\%$ regression at internal component boundaries.
\end{itemize}

This mixed performance --- averaging $\mathbf{+4.61\%}$ across the corpus --- shows that greedy
operators, such as physical reallocation, can introduce additional network cascade hops
(\texttt{CONNECTS\_TO} extensions) that backfire on isolated components even while the system-level
SRI improves. We report this result prominently, rather than only the flattering aggregate, because it
is the most direct evidence for the paper's own highest-priority future-work item: a strict per-edit
acceptance filter (\S11.2).

\subsection{Computational Overhead and CI/CD Feasibility (RQ5)}

\begin{table}[htbp]
\centering
\caption{Measured wall-clock time of the full analyze--prescribe--verify loop, single run per scenario.}
\label{tab:runtime}
\begin{tabular}{@{}lc@{}}
\toprule
Scenario & Elapsed (s) \\
\midrule
S01 Autonomous Vehicle & 13.4 \\
S02 IoT Smart City & 64.3 \\
S03 Financial Trading & 9.1 \\
S04 Healthcare & 4.7 \\
S05 Hub-and-Spoke & 18.1 \\
S06 Microservices Mesh & 12.1 \\
S07 Hyper-Scale Enterprise & 649.6 \\
\bottomrule
\end{tabular}
\end{table}

For six of the seven scenarios (S01, S03--S06), the full analysis--generation--verification loop
completes in under 20 seconds; S02 (1322 structural edges) completes in 64.3 seconds; S07 (300
processes, 3245 structural edges) completes in 649.6 seconds ($\sim$10.8 minutes). Bypassing the Neo4j
instance with the in-memory \texttt{MemoryRepository} is the enabling optimization. Measured on an
Intel Core i7-1370P (14 cores / 20 threads), 32 GB RAM, Ubuntu Linux, Python 3.11.5, single-threaded
execution.

The standalone CI/CD detection gate script (\texttt{detect\_antipatterns.py}, which runs the
twenty-one detectors of \S4 but not the full generate--verify loop) is considerably faster, since
blocking gates must survive as sub-minute checks without frustrating engineering teams:

\begin{itemize}
\item \textbf{Tiny / small scales} ($\le 25$ vertices): $< 2$ seconds.
\item \textbf{Medium scale} ($\sim$50 components, S04-like): $\approx 5$ seconds.
\item \textbf{Large scale} (80--100 components): $\approx 12$ seconds.
\item \textbf{Xlarge scale} (150--300 components, S07-like): $\approx 40$ seconds.
\end{itemize}

Taken together, these figures are compatible with pre-deployment CI/CD gating at every evaluated
scale: the detection gate itself stays under a minute even at hyper-scale, and the heavier
generate--verify loop --- run less frequently, e.g. nightly or on-demand per merge request rather than
on every commit --- fits comfortably within nightly batch windows even at 300 processes.

% ===========================================================================
\section{Discussion and Threats to Validity}
% ===========================================================================

\subsection{What Naming and Verifying Buys}

Two results, read together, motivate the paper's two-stage design. \S9.1 shows that a \emph{named,
severity-tiered} finding --- not just a scalar criticality number --- correlates strongly with
independent ground truth, giving architects a specific pathology to act on rather than an opaque
score. \S9.4 then shows why acting on a name is not enough on its own: each recommendation must also
be \emph{verified} against a cascade model before it reaches the architect, converting refactoring
recommendation from a suggestion service into a quality-evaluation instrument. An unverified
recommender, naming the correct anti-patterns but skipping verification, would have reported S05 and
S07 as unambiguous wins at the system level while silently degrading a quarter to a third of their
individual components.

\subsection{Positioning in CI/CD and Technical-Debt Workflows}

Two claim boundaries govern responsible deployment of the framework. First, \textbf{prescriptive
recommendations are advisory}: the pipeline surfaces verified refactoring blueprints for architect
review (\S6.2, Stage 7), but does not auto-apply mutations, since verified-in-simulation does not
entail correct-in-production. Second, \textbf{blocking-gate claims are reserved for anti-pattern
detection} (\S7) --- new, unwaived findings introduced relative to the merge base --- while composite
criticality rankings and prescriptive recommendations inform but do not block. The gate blocks new
architectural debt; the recommender proposes verified repayments of existing debt. For teams without
full CI/CD automation, the catalog of \S4 additionally functions as a structured architecture-review
checklist: twenty-one specific, testable questions about the system's graph structure, in the spirit
of design review by checklist as practiced in aviation and surgery, rather than an informal ``does
this look healthy?'' pass.

\subsection{Construct Validity}

Both detection and verification are defined relative to the same discrete-event cascade simulator:
detection's validation (\S4.5, \S9.1) demonstrates that named patterns correlate with \emph{simulated}
impact, and SRI reductions (\S9.2) demonstrate that prescriptions improve resilience \emph{as the
simulator models it} --- neither necessarily as a production system would experience it. We mitigate
this by grounding both the catalog's patterns and the prescriptive operators in mechanisms meaningful
independent of the simulator: established dependability practice (fan-out reduction, anti-affinity
scheduling, QoS hardening) and graph-theoretic reliability results predating this work
\cite{Baldwin2000,Colbourn1987}. Because the same simulator produces both diagnosis and verification
for the prescriptive loop, that loop is internally consistent by construction; the transfer of
verified gains to operational deployments is an external question (\S10.5).

\subsection{Internal Validity}

The prescriptive engine is rule-based and greedy: it emits every triggered mutation rather than
searching the policy space, so reported system-level gains are a lower bound on what an optimizing
search could achieve, and no optimality claim is made. Verification uses identical fault scenarios and
seeds for baseline and mutated topologies, so paired comparisons are not confounded by scenario
sampling. The independence guarantee of \S6.4 rules out circular leakage between diagnostics and
ground-truth labels. Per-type label degeneracy (passive shared libraries pooled with active
applications) can mask signals in aggregate statistics; as in \cite{SaG}, we treat stratified
reporting as mandatory wherever type-level or pattern-level results are given, and both \S4.5's
pattern-specific evidence and \S9.4's component-level breakdown are instances of this discipline
applied to detection and prescription respectively.

\subsection{External Validity}

All eight (detection) and seven (prescription) scenarios are parameterized synthetic topologies. While
the presets mimic representative domain verticals, they may not capture the runtime complexity of
industrial clusters --- dynamic workload shifts, packet-loss bursts, transient hardware faults. The
catalog itself is scoped to the publish--subscribe communication paradigm: systems combining pub-sub
with request-response patterns (hybrid microservices, mixed REST/event architectures) will require
additional patterns addressing the request-response side, and \texttt{QOS\_MISMATCH} is currently
specified for DDS/ROS~2 and MQTT QoS-weight semantics, requiring adaptation for other middleware
platforms. Validating prescriptions against a real system (e.g. replaying the engine on the ATM
topology reserved for external validation in \cite{SaG}) is the most direct extension.

\subsection{Conclusion Validity}

The prescriptive-evaluation scenario sample is small ($n=7$); the Wilcoxon result reported in \S9.2
($W=0.0$, $p=0.0156$) is the smallest attainable two-sided $p$-value at this sample size, since all
seven deltas share the same sign --- it should be read as consistent directional evidence rather than
as a claim of large effect size. Results are single-run per scenario under a verified-deterministic
simulator configuration ($\sigma_{\text{seed}}=0$, \S8.3); this determinism is a property of the
current cascade-probability-1.0 configuration and would need re-establishing under the stochastic
propagation model proposed as future work (\S11.2).

\subsection{Engineering Trade-offs}

The closed-loop verification step adds computational cost over emitting unverified recommendations.
Since the framework targets pre-deployment optimization rather than runtime alerting, \S9.5 shows this
cost stays within CI/CD budgets at every evaluated scale, and we consider verification the feature
that distinguishes a prescription from a suggestion --- just as naming a finding against an empirically
validated catalog is the feature that distinguishes a diagnosis from an unlabeled score.

% ===========================================================================
\section{Conclusion and Future Work}
% ===========================================================================

\subsection{Conclusions}

This paper presented SaG-Prescribe, a graph-based framework that unifies the detection of named,
severity-tiered architectural anti-patterns with closed-loop, DevOps-integrated prescriptive
refactoring for distributed publish--subscribe architectures. The twenty-one-pattern catalog gives
practitioners a pub-sub-specific vocabulary analogous to established object-oriented and microservices
smell catalogs, with detection rules empirically validated against failure-simulation ground truth
(Spearman $\rho = 0.876$ overall, $0.943$ at scale; $F_1 = 0.923$). By augmenting these topological
diagnostics with a code-quality bridge (CQP) and extending the Software-as-a-Graph baseline, we
introduced an end-to-end pipeline that translates named findings into three concrete mutation operators
--- topic splitting, anti-affinity reallocation, and QoS hardening --- verifies every compiled policy
against the same discrete-event cascade oracle that produced the diagnosis, and operationalizes the
result as a delta-aware CI/CD quality gate. Across seven benchmark scenarios, the generated
prescriptions reduce the System Risk Index in every case (1.4\%--15.9\% relative, Wilcoxon $W=0.0$,
$p=0.0156$) within runtime envelopes compatible with merge-request-gated or nightly CI/CD pipelines,
while an honest component-level audit (\S9.4) shows the current greedy policy is not yet uniformly
beneficial below the system level.

\subsection{Future Work}

\begin{enumerate}
\item \textbf{Per-edit acceptance filter (highest priority):} eliminating counterproductive mutations
  before policy execution, directly motivated by the component-level regressions of \S9.4.
\item \textbf{Budget-constrained policy search:} incorporating explicit cost models and modification
  budgets ($\mathcal{B} < \infty$ in \S5), turning the engine from exhaustive rule firing into subset
  selection over candidate mutations.
\item \textbf{Stochastic cascade propagation:} the current simulator is deterministic under its
  default configuration (cascade probability 1.0), so seed variance is trivially zero (\S8.3).
  Introducing probabilistic propagation ($<1.0$) would make the $\kappa\cdot\sigma_{\text{seed}}$
  acceptance criterion of \S5 load-bearing and let RQ2 test robustness to genuine simulation noise.
\item \textbf{Operator ablation and learned policy ordering:} isolating per-operator $\Delta$SRI
  contributions and exploring learned prioritization over the rule-based candidate set.
\item \textbf{Catalog extension:} generalizing the pattern set to hybrid REST/event architectures and
  to middleware platforms whose QoS semantics differ from the DDS/ROS~2/MQTT weight formula
  \texttt{QOS\_MISMATCH} currently assumes (\S10.5).
\item \textbf{Real-system replication:} applying the engine to the ATM topology of \cite{SaG} and to
  harvested industrial configurations, closing the external-validity gap of \S10.5.
\item \textbf{LLM-assisted pull-request generation:} linking large language model code assistants to
  automatically generate pull requests that implement the synthesized architectural refactoring
  blueprints. This is the paper's only LLM-adjacent contribution and should not be oversold beyond
  this single future-work item --- the paper makes no LLM or AutoML contribution today.
\end{enumerate}

% ===========================================================================
\begin{thebibliography}{99}
% ===========================================================================

\bibitem{SaG} [Authors]. \emph{Software-as-a-Graph: A Static System Analysis Framework for
  Pre-Deployment Quality Gating and Failure Simulation of Publish-Subscribe Middleware.} Journal of
  Systems and Software, under review / to appear. [Update status at submission time; AuSE permits
  citing companion work under review with a copy supplied to the editor. Confirm submission status
  and disclose per AuSE's originality/overlap policy in the cover letter.]

\bibitem{Harman2012} M. Harman, S. A. Mansouri, Y. Zhang, ``Search-Based Software Engineering: Trends,
  Techniques and Applications,'' \emph{ACM Computing Surveys}, 45(1), Article 11, 2012.

\bibitem{Aleti2013} A. Aleti, B. Buhnova, L. Grunske, A. Koziolek, I. Meedeniya, ``Software
  Architecture Optimization Methods: A Systematic Literature Review,'' \emph{IEEE Transactions on
  Software Engineering}, 39(5), 658--683, 2013.

\bibitem{Pallickara2007} S. Pallickara, H. Bulut, G. Fox, ``Fault-Tolerant Reliable Delivery of
  Messages in Distributed Publish/Subscribe Systems,'' \emph{Proc. 4th IEEE International Conference
  on Autonomic Computing (ICAC 2007)}, 2007.

\bibitem{Chang2014} T. Chang, S. Duan, H. Meling, S. Peisert, H. Zhang, ``P2S: A Fault-Tolerant
  Publish/Subscribe Infrastructure,'' \emph{Proc. 8th ACM International Conference on Distributed
  Event-Based Systems (DEBS 2014)}, 2014.

\bibitem{Wang2015} G. Wang, J. Koshy, S. Subramanian, K. Paramasivam, M. Zadeh, N. Narkhede, J. Rao,
  J. Kreps, J. Stein, ``Building a Replicated Logging System with Apache Kafka,'' \emph{Proceedings of
  the VLDB Endowment}, 8(12), 1654--1655, 2015.

\bibitem{Lee2025Latency} S. Lee, H.-S. Park, J. Chae, K.-J. Park, ``Probabilistic Latency Analysis of
  the Data Distribution Service in ROS 2,'' \emph{arXiv:2508.10413}, 2025.

\bibitem{Lee2025QoS} S. Lee, J. Kang, K.-J. Park, ``Dependency Chain Analysis of ROS 2 DDS QoS
  Policies: From Lifecycle Tutorial to Static Verification,'' \emph{arXiv:2509.03381}, 2025.

\bibitem{Freeman1977} L. C. Freeman, ``A set of measures of centrality based on betweenness,''
  \emph{Sociometry}, vol. 40, no. 1, pp. 35--41, 1977.

\bibitem{Bakhtin2025} A. Bakhtin, M. Esposito, V. Lenarduzzi, D. Taibi, ``Network Centrality as a New
  Perspective on Microservice Architecture,'' \emph{Proc. IEEE International Conference on Software
  Architecture (ICSA 2025)}, 72--83, 2025.

\bibitem{Falci2017} D. H. M. Falci, O. A. Gomes, F. S. Parreiras, ``Complex Networks Analysis for
  Software Architecture: an Hibernate Call Graph Study,'' \emph{arXiv:1706.09859}, 2017.

\bibitem{Saaty1980} T. L. Saaty, \emph{The Analytic Hierarchy Process: Planning, Priority Setting,
  Resource Allocation}, McGraw-Hill, 1980.

\bibitem{Fowler1999} M. Fowler, \emph{Refactoring: Improving the Design of Existing Code},
  Addison-Wesley, 1999.

\bibitem{Brown1998} W. H. Brown, R. C. Malveau, H. W. McCormick, T. J. Mowbray, \emph{AntiPatterns:
  Refactoring Software, Architectures, and Projects in Crisis}, Wiley, 1998.

\bibitem{Suryanarayana2014} G. Suryanarayana, G. Samarthyam, T. Sharma, \emph{Refactoring for Software
  Design Smells: Managing Technical Debt}, Morgan Kaufmann, 2014.

\bibitem{Richardson2018} C. Richardson, \emph{Microservices Patterns: With Examples in Java}, Manning,
  2018.

\bibitem{Taibi2020} D. Taibi, V. Lenarduzzi, C. Pahl, ``Microservices anti-patterns: A taxonomy,'' in
  \emph{Microservices: Science and Engineering}, Springer, 2020.

\bibitem{Baldwin2000} C. Y. Baldwin, K. B. Clark, \emph{Design Rules, Volume 1: The Power of
  Modularity}, MIT Press, 2000.

\bibitem{Colbourn1987} C. J. Colbourn, \emph{The Combinatorics of Network Reliability}, Oxford
  University Press, 1987.

\bibitem{Lehman1996} M. M. Lehman, ``Laws of software evolution revisited,'' \emph{Proceedings of
  EWSPT '96}, Springer, 1996.

\bibitem{Martin2003} R. C. Martin, \emph{Agile Software Development, Principles, Patterns, and
  Practices}, Prentice Hall, 2003.

\bibitem{Nygard2018} M. T. Nygard, \emph{Release It! Design and Deploy Production-Ready Software}
  (2nd ed.), Pragmatic Bookshelf, 2018.

\end{thebibliography}

\begin{quote}
\emph{[Reference-list note, not part of the manuscript: references [2]--[12] were sourced from the
previously verified prescriptive/SBSE/pub-sub bibliography; references [13]--[22] are sourced from
\texttt{docs/antipatterns.md}'s own verified bibliography and are real, non-invented candidates, though
full text was not re-read for each in this pass --- sanity-check relevance before submission. AuSE
reviewers will expect $\sim$30--45 references total; the two [REF: \ldots] placeholders in \S2.3
(learning-based and LLM-based refactoring recommenders) still need populating with real citations, and
this list should otherwise be expanded further before submission.]}
\end{quote}

\section*{Declarations}

\begin{itemize}
\item \textbf{Funding:} [to be completed]
\item \textbf{Competing interests:} [to be completed]
\item \textbf{Data availability:} A replication package containing scenario configurations, seeds, and
  the prescriptive pipeline implementation will be made available at [URL pending].
\item \textbf{Ethics approval:} Not applicable.
\end{itemize}

\end{document}
